\section{Introduction}
\label{sec:intro}

% P1. Percussive sound and why sinusoids plus noise fall short.
Percussive sounds are produced by striking vibrating solids such as membranes, bars, plates, and shells, and their timbre varies widely with the shape, size, and material of the struck object \citep{fletcher1998physics}.
Because the instrument vibrates freely after the strike, these sounds share a decaying envelope. 
Modeling and synthesis of percussive audio date back to analog drum machines such as the Roland TR-808 \citep{werner2014physically,werner2014tr808} and have been approached through physical modeling \citep{bilbao2010percussion}, signal processing \citep{karplus1983digital}, and, most recently, neural audio synthesis (NAS), in which deep generative models learn to synthesize drum sounds from data \citep{aouameur2019neural,drysdale2020adversarial,nistal2020drumgan,rouard2021crash}.
Spectral modeling synthesis (SMS) \citep{serra1990spectral} represents a sound as sinusoidal partials with slowly varying amplitudes and frequencies \citep{mcaulay1986speech} plus a stochastic residual, and differentiable digital signal processing (DDSP) \citep{engel2020ddsp} implements a synthesizer of this kind with differentiable operations, so that a neural network can be trained end to end to predict its parameters from audio while each synthesis module keeps an interpretable role.
Applied to percussion, this formulation represents the modes of the struck body as sinusoids and the remainder as stochastic noise, which suits instruments with a clear pitch and prominent modes, such as membranophones.
For tom-toms, a mode-wise analysis and synthesis based on the discrete cosine transform has even produced samples that listeners could not distinguish from recordings \citep{kirby2021evolution}.
Idiophones such as cymbals, in contrast, consist of hundreds or even thousands of modes that give them a noisy high-frequency range, and their sound blooms as the vibration moves from linear to nonlinear regimes \citep{skare2019realtime,chaigne2005nonlinear}, which makes them difficult to analyze.
The strike itself also produces a short broadband burst, the \textit{transient}, so a percussive sound combines narrow-band components with components that change rapidly over time, and neither is easily captured by sinusoids and noise alone \citep{verma2000extending,barahonarios2024noisebandnet}.

% P2. Differentiable responses to these difficulties and the remaining problem.
DDSP has since been adopted in a wide range of synthesis approaches \citep{hayes2024review,shan2022differentiable,caspe2022ddx7,barahonarios2024noisebandnet}, including several that extend it to inharmonic sounds \citep{renault2023ddsp,diaz2023rigid,hayes2023sinusoidal}.
Transients have likewise been given a dedicated component, both in classical sines, transients, and noise representations \citep{levine1998sines,verma2000extending} and in differentiable models \citep{liu2024ddspsfx,simionato2025sines}.
DrumBlender \citep{shier2023differentiable} brought this sines, noise, and transient decomposition to percussion, yet cymbals remained the hardest class to reconstruct.
With a limited set of sinusoidal tracks, much of the dense modal content of a cymbal is left for the stochastic noise component, so a single model that reconstructs both pitched drums and cymbals faithfully remains difficult.

% P3. Timbre control and the gap.
Reconstruction alone, however, does not give a sound designer a way to change a sound.
Percussion synthesis has offered timbre control through the fixed circuits of analog drum machines \citep{werner2014physically,werner2014tr808} and through timbral descriptors and latent spaces of neural drum generators \citep{aouameur2019neural,ramires2020neural,nistal2020drumgan,lavault2022stylewavegan}.
Outside percussion, control has also been obtained from latent spaces learned without timbre labels \citep{caillon2021rave,kotowski2026plaud,paek2025learning}.
An analysis--synthesis model such as DrumBlender \citep{shier2023differentiable} reconstructs a given sample through interpretable parameters, but these parameters are low-level and numerous, and high-level control was left as future work.
Turning these parameters into high-level timbre controls requires controls that are learned without labels, start from the reconstruction of the given sample, and act through the parameters of the synthesizer.

% P4. This work.
To this end, this paper builds on DrumBlender \citep{shier2023differentiable} and reconstructs a percussive one-shot as the sum of modal, noise, and transient components, which are trained jointly through their summed output.
The modal component starts from sinusoidal tracks obtained by a hybrid of constant-Q transform (CQT) and short-time Fourier transform (STFT) analysis, and a neural network predicts bounded corrections to their frequencies and amplitudes.
The noise component is a bank of narrow noise bands following NoiseBandNet \citep{barahonarios2024noisebandnet}, whose amplitudes are predicted with the corrected modal parameters as an additional input, and a small temporal convolutional network adds a transient computed from the modal waveform, which makes the model a hybrid of signal processing and neural synthesis.
After reconstruction training, a sparse autoencoder (SAE) is fitted to the shared representation of the frozen model, and its temporal features define continuous edits around the reconstruction of each sample.
These edits are evaluated on held-out one-shots, together with the separate effects of the modal-to-noise and transient conditioning paths.
Trained on about 17,000 one-shots collected from public datasets and commercial sample packs, the model \textbf{[TBD: reconstruction result relative to DrumBlender]}, and its SAE features provide controls that \textbf{[TBD: control result]}.

The contributions of this paper are as follows.
\begin{itemize}
  \item A hybrid modal analysis front-end that tracks both low-frequency modes and dense, short-lived high-frequency modes (\sectionref{sec:method-analysis}).
  \item A jointly trained modal, noise, and transient synthesizer in which bounded corrections refine the analyzed tracks and the noise and transient components are conditioned on the corrected modal component to model unlabeled diverse percussive audio (\sectionref{sec:method-synthesis}).
  \item Sample-anchored timbre controls from a post-hoc temporal SAE, with separate interventions to measure how a latent edit propagates through modal conditioning of noise and the transient clip code (\sectionref{sec:method-control}).
\end{itemize}
